% Statement of interest -- Xanadu Residency Program (deadline Nov 15, 2026, 11:59pm ET).
% Build: cd report && tectonic statement_of_interest.tex
\documentclass[11pt,a4paper]{article}
\usepackage[hmargin=2cm,vmargin=1.8cm]{geometry}
\usepackage[T1]{fontenc}
\usepackage{lmodern}
\usepackage{microtype}
\usepackage{xcolor}
\usepackage[colorlinks=true,urlcolor=blue!50!black]{hyperref}
\setlength{\parindent}{0pt}
\setlength{\parskip}{5pt}
\newcommand{\fillin}[1]{\textcolor{orange!80!black}{[#1]}}
\pagestyle{empty}

\begin{document}

{\large\textbf{Statement of interest -- Xanadu Residency Program}}\\
Tuan Bui \;\textbullet\; Quantum Technology and Computer Science, Aalto University (2024--2027)\\
\href{mailto:tuan.bui@aalto.fi}{tuan.bui@aalto.fi} \;\textbullet\; \fillin{GitHub link to the project}

\textbf{Area and topic.} Research: \emph{characterising the differences between quantum and
classical machine learning models}. I would most like to work with the quantum machine learning
team, ideally with Joseph Bowles, whose paper with Shahnawaz Ahmed and Maria Schuld,
\emph{Better than classical? The subtle art of benchmarking quantum machine learning models}, is
the direct starting point of the project described below. I would also be glad to contribute
PennyLane functionality or demos that come out of this work.

\textbf{What I have done.} To learn how to ask ``where might a quantum model help?'' carefully,
I built a fair regression benchmark on four real datasets chosen along one axis: how periodic
the inputs are. These are tidal water levels at San Francisco (inputs are astronomical phases),
wind-turbine power (wind speed and direction), expected goals in soccer (shot geometry), and
galaxy redshift from photometry (no periodic inputs, used as a control). I compared a data
re-uploading circuit and a fidelity quantum kernel, both in PennyLane, with five classical models
and with the methods each field actually uses: UTide harmonic analysis, IEC-style power
curves, logistic xG and $k$NN photo-$z$. Every model had the same Optuna budget, cross-validation
respected time blocks and matches, and periodic inputs were encoded equivalently for both
families: 850 runs in total, with 50 to 1000 training points.

\textbf{What I found.} No quantum model beat the best classical or domain baseline anywhere.
The more interesting results are about \emph{why}:
\begin{itemize}\setlength{\itemsep}{0pt}\setlength{\topsep}{2pt}
  \item \textbf{The feature map has to match the data.} The quantum kernel with an angle map is
        competitive on tides (test $R^2$ 0.927 vs.\ 0.931 for a Gaussian process). The IQP map
        failed on periodic inputs: its $ZZ(x_jx_k)$ terms are not $2\pi$-periodic. The tuner
        rejected IQP in every tides and wind run, but often chose it for the non-periodic photo-$z$ data.
  \item \textbf{Expressivity vs.\ trainability.} The variational circuit stalled on tides. Using
        the Fourier view of re-uploading models (Schuld, Sweke \& Meyer), I computed
        least-squares ``spectrum ceilings'' and exact model spectra by FFT. A one-layer
        circuit's frequencies already allow $R^2 = 0.932$, so the gap was not missing
        frequencies. It came from training (a longer budget raised the circuit's mean $R^2$ from 0.870
        to 0.894) and from the limited set of Fourier coefficients the ansatz can reach. Deeper circuits put about 30\,\% of their output at frequencies the
        data does not contain, and trained worse.
  \item My first hypothesis (that tides need higher harmonics) was wrong; the analysis
        showed it. I value that kind of result as much as a positive one.
\end{itemize}

\textbf{What I would like to do at Xanadu.} I would like to turn this kind of analysis into
reusable tools and sharper questions. Examples: spectrum-ceiling diagnostics that tell a user
in advance whether a circuit's frequencies suit their data; feature maps designed to respect
periodic and mixed-type inputs; and the step from noiseless simulation to shot noise and
hardware-realistic budgets. I am equally happy to work on the team's current questions; I care
more about getting the comparison right than about the answer favouring quantum.

\textbf{Why me.} My coursework at Aalto covers quantum mechanics, quantum information, practical
quantum computing and deep learning. As a research assistant intern at the Quantum AI Lab in Hanoi (2025) I worked with the
quantum optimal control and superconducting-qubit teams. In a second project I raised a 6-qubit
variational phishing classifier's accuracy from 77\% to 89\% by tracing its shortfall to the SPSA
optimizer rather than the circuit, and showed with density-matrix simulation that depolarizing
noise shifts its decision threshold while leaving its ranking (AUROC) almost unchanged. The habit
behind both projects, finding out \emph{why} a quantum model does or does not work, is what I
would bring to the team. I have been a teaching assistant for Aalto courses in algorithms and in
statistical physics, and at the Vietnam Quantum Summer School. I write tested, reproducible
Python: the benchmark's code, 50 unit tests, and a paper whose numbers are all generated from the
results are at \fillin{repository link}.

\end{document}
